\documentclass[10pt,letterpaper]{article}
\usepackage{cornell-notes}

\title{Clause 30.04.02.02.02: Ctype Members}
\author{}
\date{}
\setDocTitle{Clause 30.04.02.02.02: Ctype Members}
\setDocSubtitle{C++ 2024}
\setDocOwner{C++ 2024 Cornell Notes}
\setDocDate{\today}
\setCornellCollection{C++ 2024 Cornell Notes}
\setCornellUnitType{Clause}
\setCornellUnitNumber{30.04.02.02.02}
\setCornellUnitTitle{Ctype Members}
\hypersetup{pdftitle={Clause 30.04.02.02.02: Ctype Members},pdfsubject={Cornell notes},pdfauthor={}}

\begin{document}
\maketitle
\begin{CornellOverviewBox}{Overview}
\textbf{Classification:} Standard library - Normative\\[2pt]
\textbf{Learning objective:} Understand the rules, terminology, and semantic consequences associated with ctype members.
\end{CornellOverviewBox}\begin{CornellSummaryBox}{Summary}
{\sffamily\bfseries\color{Accent} Summary}\par\vspace{3pt}Section 30.4.2.2.2 focuses on ctype members. Apply the specified interface together with its mandates, effects, result conditions, exception behavior, and complexity constraints. When applying it, preserve the distinction between mandatory requirements, recommendations, permissions, and behavior deliberately left undefined, unspecified, or implementation-defined.
\end{CornellSummaryBox}

\end{document}
